\renewcommand{\tablename}{Quadro}
\captionsetup{labelsep=endash, labelfont=bf}

\begin{center}
\renewcommand{\arraystretch}{1.3} % Reduzido para caber melhor

\begin{longtable}{|p{5cm}|*{13}{>{\centering\arraybackslash}p{0.48cm}|}}
\caption{Cronograma de atividades do TCC}
\label{tab:cronograma-rotacionado} \\
\hline
\rowcolor{gray!30}
\textbf{Etapas} & 
\rotatebox{90}{\textbf{Mar/25}} & 
\rotatebox{90}{\textbf{Abr/25}} & 
\rotatebox{90}{\textbf{Mai/25}} & 
\rotatebox{90}{\textbf{Jun/25}} & 
\rotatebox{90}{\textbf{Jul/25}} & 
\rotatebox{90}{\textbf{Ago/25}} & 
\rotatebox{90}{\textbf{Set/25}} & 
\rotatebox{90}{\textbf{Out/25}} & 
\rotatebox{90}{\textbf{Nov/25}} & 
\rotatebox{90}{\textbf{Dez/25}} & 
\rotatebox{90}{\textbf{Jan/26}} & 
\rotatebox{90}{\textbf{Fev/26}} & 
\rotatebox{90}{\textbf{Mar/26}} \\
\hline
\endfirsthead

% Corpo da tabela
\hline
Escolha do tema & X & & & & & & & & & & & & \\
\hline
Levantamento bibliográfico & X & X & & & & & & & & & & & \\
\hline
Definição do problema e objetivos & & X & X & & & & & & & & & & \\
\hline
Introdução & & & X & X & & & & & & & & & \\
\hline
Fundamentação Teorica & & & & X & X & & & & & & & & \\
\hline
Metodologia & & & & & X & X & X & X & & & & & \\
\hline
Foirmulario & & & & & & & & X & X & & & & \\
\hline
Conclusão & & & & & & & & & & X & X & & \\
\hline
Revisão geral do trabalho & & & & & & & & & & & & X & X  \\
\hline
Preparação da apresentação/defesa & & & & & & & & & & & & X & X \\
\hline
\end{longtable}

\vspace{0.2cm}
\small
\textit{Fonte: Elaborado pelo autor, 2025.}
\end{center}

